\chapter{Potential Outcomes and Estimands}
\label{chap:potential-outcomes}

\section{The framework}
\label{sec:po-framework}

Let $i = 1, \dots, n$ index units drawn independently from a population. Each
unit has a binary treatment $A_i \in \{0,1\}$, a vector of pre treatment
covariates $X_i \in \mathcal{X}$, and an observed outcome $Y_i \in \mathbb{R}$.
The framework of this book posits, in addition, a pair of potential
outcomes\index{potential outcomes} $\left( Y_i(0), Y_i(1) \right)$ representing
the outcome that would be observed under each treatment value.

\begin{assumption}[Consistency]
\label{ass:consistency}
\index{consistency assumption}
The observed outcome equals the potential outcome under the received treatment,
$Y_i = A_i Y_i(1) + (1 - A_i) Y_i(0)$.
\end{assumption}

Assumption~\ref{ass:consistency} looks like a definition and is not. It asserts
that the treatment is well enough defined that a single potential outcome
attaches to each value of $A_i$. A treatment such as ``obesity''\index{treatment version} fails this
test, because the outcome under a reduction in body mass depends on how the
reduction was achieved, and the framework cannot distinguish those routes.

\begin{assumption}[No interference]
\label{ass:sutva}
\index{interference}
The potential outcomes of unit $i$ do not depend on the treatments assigned to
other units.
\end{assumption}

Assumption~\ref{ass:sutva} is violated whenever treatment spreads. Vaccination,
job training in a local labour market, and any intervention delivered to a
network all violate it, and Section~\ref{sec:design-interference} sketches what
can be recovered when it fails.

\section{Estimands}
\label{sec:po-estimands}

\begin{definition}[Average treatment effect]
\label{def:ate}
\index{average treatment effect}
The average treatment effect is
\begin{equation*}
  \ATE = \E\left[ Y(1) - Y(0) \right].
\end{equation*}
\end{definition}

\begin{definition}[Effect on the treated]
\label{def:att}
\index{average treatment effect!on the treated}
The average treatment effect on the treated is
\begin{equation*}
  \ATT = \E\left[ Y(1) - Y(0) \mid A = 1 \right].
\end{equation*}
\end{definition}

The two differ whenever treatment assignment correlates with the effect rather
than merely with the level of the outcome, which is the usual situation when
people select into treatment on the basis of expected benefit. A policy question
about expanding a programme to those not currently receiving it is a question
about neither estimand exactly, and Section~\ref{sec:design-late} returns to
this point.

\begin{remark}
\label{rem:estimand-first}
Choosing the estimand before choosing the estimator is not a stylistic
preference. A propensity weighted difference in means and an outcome regression
fitted to the same data estimate different quantities under misspecification,
and reporting whichever is more favourable is a form of specification search
that no standard error accounts for.
\end{remark}

\section{Identification under ignorability}
\label{sec:po-identification}

The fundamental problem of causal inference\index{fundamental problem of causal inference} is that $Y_i(1)$ and $Y_i(0)$ are never both observed.
Identification proceeds by assuming that, within strata of covariates, the
treated and untreated are exchangeable.

\begin{assumption}[Conditional ignorability]
\label{ass:ignorability}
\index{ignorability}
$\left( Y(0), Y(1) \right) \indep A \mid X$.
\end{assumption}

\begin{assumption}[Positivity]
\label{ass:positivity}
\index{positivity}
There exists $\delta > 0$ such that
$\delta < \Prob(A = 1 \mid X = x) < 1 - \delta$ for all $x$ in the support of
$X$.
\end{assumption}

\begin{theorem}[Identification of the ATE]
\label{thm:ate-identification}
Under Assumptions~\ref{ass:consistency}, \ref{ass:sutva},
\ref{ass:ignorability}, and~\ref{ass:positivity},
\begin{equation}
  \ATE = \E\Big[ \E\left[ Y \mid A = 1, X \right]
               - \E\left[ Y \mid A = 0, X \right] \Big],
  \label{eq:g-formula}
\end{equation}
and the right hand side is a functional of the observed data distribution alone.
\end{theorem}

\begin{proof}
Fix $a \in \{0,1\}$. By the tower property,
$\E[Y(a)] = \E\left[ \E\left[ Y(a) \mid X \right] \right]$. By
Assumption~\ref{ass:ignorability}, $Y(a) \indep A \mid X$, so
$\E\left[ Y(a) \mid X \right] = \E\left[ Y(a) \mid A = a, X \right]$. The
conditioning event has positive probability by
Assumption~\ref{ass:positivity}. By Assumption~\ref{ass:consistency}, on the
event $\{A = a\}$ we have $Y = Y(a)$, hence
$\E\left[ Y(a) \mid A = a, X \right] = \E\left[ Y \mid A = a, X \right]$.
Taking the difference at $a = 1$ and $a = 0$ and applying linearity of
expectation gives~\eqref{eq:g-formula}. Every term on the right involves only
the joint law of $(Y, A, X)$.
\end{proof}

The right hand side of~\eqref{eq:g-formula} is the g formula\index{g formula}.
Its structure is worth noting: the inner difference is a regression contrast at
fixed $X$, and the outer expectation is over the marginal covariate
distribution, not over the covariate distribution within either treatment arm.
Standardising to the wrong marginal is the most common error in applied practice
and produces a quantity that is neither the ATE nor the ATT.

\begin{corollary}[Identification of the ATT]
\label{cor:att-identification}
Under the same assumptions, with Assumption~\ref{ass:positivity} weakened to
$\Prob(A = 1 \mid X = x) < 1 - \delta$,
\begin{equation*}
  \ATT = \E\Big[ \E\left[ Y \mid A = 1, X \right]
               - \E\left[ Y \mid A = 0, X \right] \,\Big|\, A = 1 \Big].
\end{equation*}
\end{corollary}

\begin{proof}
Repeat the proof of Theorem~\ref{thm:ate-identification} with every expectation
conditioned on $A = 1$. Only $\E[Y(0) \mid A = 1]$ requires the ignorability
assumption, and only the upper positivity bound is used, because units with
$\Prob(A = 1 \mid X) = 0$ contribute no mass to the conditioning event.
\end{proof}

\section{What the assumptions cost}
\label{sec:po-cost}

Assumption~\ref{ass:ignorability} is untestable. No function of the observed
data distinguishes a world in which it holds from one in which it fails, because
the failure concerns the joint law of quantities that are never jointly
observed. This is not a technicality to be noted and moved past. It is the
reason Part~\ref{part:design} exists.

Assumption~\ref{ass:positivity} is testable in principle and routinely violated
in practice \citep{ferreiro2012positivity}. Table~\ref{tab:overlap} reports the overlap diagnostics for the
worked example carried through this book, a study of a regional retraining
programme in which enrolment was neither randomised nor administratively
rationed.

\begin{table}[t]
\centering
\caption[Overlap diagnostics for the retraining study]{Overlap diagnostics for
the retraining study, $n = 8412$. The estimated propensity score is binned into
deciles. Positivity is marginal in the top decile, where 96 percent of units are
treated, and this is the stratum that drives the instability of the weighted
estimators in Chapter~\ref{chap:propensity}.}
\label{tab:overlap}
\begin{tabular}{lrrrr}
\toprule
Decile of $\hat{e}(X)$ & Range & Treated & Control & Treated share \\
\midrule
1  & 0.01 to 0.06 & 22  & 819 & 0.026 \\
2  & 0.06 to 0.11 & 61  & 780 & 0.073 \\
3  & 0.11 to 0.18 & 118 & 723 & 0.140 \\
4  & 0.18 to 0.26 & 190 & 651 & 0.226 \\
5  & 0.26 to 0.35 & 264 & 577 & 0.314 \\
6  & 0.35 to 0.46 & 351 & 490 & 0.417 \\
7  & 0.46 to 0.59 & 452 & 390 & 0.537 \\
8  & 0.59 to 0.73 & 559 & 282 & 0.665 \\
9  & 0.73 to 0.88 & 682 & 159 & 0.811 \\
10 & 0.88 to 0.99 & 808 &  33 & 0.961 \\
\midrule
Total & & 3507 & 4904 & 0.417 \\
\bottomrule
\end{tabular}
\end{table}

Two features of Table~\ref{tab:overlap} matter later. The support is nominally
complete, in that both arms are represented in every decile. The effective
support is not: the top decile contributes 33 control units to an estimate that
must extrapolate across 808 treated units, and any estimator that weights by
$1/(1 - \hat{e})$ will place enormous weight on those 33 observations. This is
the mechanism behind the variance inflation documented in
Section~\ref{sec:ps-variance}.
